\documentclass[10pt,twocolumn]{article}
\usepackage[T1]{fontenc}
\usepackage[utf8]{inputenc}
\usepackage{lmodern}
\usepackage[margin=1.8cm,top=2.4cm,bottom=2cm,columnsep=0.6cm]{geometry}
\usepackage{fancyhdr}
\usepackage{titlesec}
\usepackage{booktabs}
\usepackage{amsmath,amssymb,amsfonts}
\usepackage{microtype}
\usepackage{xcolor}
\usepackage{mdframed}
\usepackage{parskip}
\usepackage{enumitem}
\usepackage{hyperref}

\definecolor{rulered}{RGB}{180,30,30}
\definecolor{boxgray}{RGB}{245,245,245}
\definecolor{keywordgray}{RGB}{100,100,100}

\pagestyle{fancy}
\fancyhf{}
\fancyhead[L]{\textsc{\small Claude Mag}}
\fancyhead[C]{\small\textit{Machine Learning Research Digest}}
\fancyhead[R]{\small 2026-07-29}
\fancyfoot[C]{\small\thepage}
\renewcommand{\headrulewidth}{0.8pt}

\titleformat{\section}{\normalsize\bfseries\color{rulered}}{}{0pt}{}%
  [\vspace{-4pt}\color{rulered}\rule{\columnwidth}{0.4pt}]
\titlespacing{\section}{0pt}{8pt}{4pt}

\hypersetup{colorlinks=true,urlcolor=rulered,linkcolor=black}

\begin{document}
\setlength{\parindent}{0pt}
\setlength{\parskip}{4pt}

{\small\color{keywordgray}\textbf{Keywords:}
  multimodal reasoning \textbullet{}
  reinforcement learning \textbullet{}
  vision-language models \textbullet{}
  geometry}\par\vspace{4pt}

{\LARGE\bfseries\raggedright
  MIRROR: Learning from the Other View
  for Multi-Modal Reasoning}\par\vspace{4pt}

{\large\itshape\raggedright
  Modality-Informed Reciprocal Reasoning Optimization Trains
  Vision-Language Models to Transfer Strengths Across
  Text, Diagram, and Combined Views of the Same Problem}\par\vspace{6pt}

{\small\textbf{By} Wen Ye, Yuxiao Qu, Aviral Kumar, Xuezhe Ma
  (University of Southern California; Carnegie Mellon University)}\par\vspace{2pt}

{\small\color{keywordgray}arXiv:2607.21552 \textbullet{} July 23, 2026}
\par\vspace{2pt}\noindent{\color{rulered}\rule{\columnwidth}{0.6pt}}\vspace{6pt}

Vision-language models fail at visual reasoning even on geometry problems
that have mathematically equivalent text, diagram, and combined views.
MIRROR, from researchers at USC and CMU, exploits this inconsistency:
rather than treating it as noise, the method identifies which view a model
currently solves best and uses that signal to lift performance on the
harder view through reciprocal cross-modal reward shaping during RL
training. On ODA-Data, a new paired multimodal geometry benchmark, MIRROR
improves image-dominant reasoning by 17 points over a strong baseline.

\begin{mdframed}[backgroundcolor=boxgray,linecolor=rulered,linewidth=1pt,
    innerleftmargin=6pt,innerrightmargin=6pt,
    innertopmargin=6pt,innerbottommargin=6pt]
\textbf{\small AT A GLANCE}\par\vspace{4pt}
\begin{description}[leftmargin=0pt,labelsep=4pt,itemsep=2pt,parsep=0pt,
    font=\small\bfseries]
  \item[Problem:] VLMs solve the same geometry problem in text but
    fail on the diagram, revealing complementary reasoning gaps across
    modalities that standard post-training ignores.
  \item[Why it matters:] Closing modality gaps enables VLMs to reason
    reliably regardless of how a problem is presented---critical for
    real-world deployment where input format is unpredictable.
  \item[Method:] MIRROR evaluates each view per-problem, identifies
    the best-performing view, and shapes rewards on weaker views
    toward the best-view distribution via RL.
  \item[Result:] +17pp on image-dominant and +9pp on combined view
    vs.\ GRPO baseline on ODA-Data geometry; consistent gains on
    three external geometry benchmarks.
  \item[Evidence:] Ablating cross-modal shaping loses 9.4pp image
    accuracy; scrambled skills confirm content drives improvement.
\end{description}
\end{mdframed}

\section{The Big Picture}

Vision-language models are post-trained on massive image-text pairs
and achieve strong results on captioning, VQA, and chart understanding.
Yet on mathematical visual reasoning---where the model must interpret
a geometric diagram and apply deductive inference---performance lags
substantially behind text-only reasoning.

The standard explanation is that vision encoding is weaker than language
encoding. MIRROR's key insight is that this framing misses a more
specific phenomenon: the \emph{same problem} in a different modality
often elicits dramatically different behavior from the same model.
A model that easily solves a geometry problem from a text description
may completely fail on the same problem presented as a diagram---not
because it lacks geometric knowledge, but because it fails to activate
that knowledge from visual input.

\section{Why Existing Approaches Fall Short}

Standard VLM post-training mixes image-text pairs and optimizes a
single objective (SFT or RLHF) over all examples, treating modality
as a property of the input rather than an independent factor in
reasoning quality. This approach:

\begin{itemize}[itemsep=2pt,parsep=0pt]
  \item \textbf{Averages over view-dependent performance:} A model
    that excels on text but fails on diagrams receives the same
    training gradient regardless of which view it was shown.
  \item \textbf{Does not exploit complementary signals:} The model's
    success on text implicitly encodes geometric reasoning that could
    guide improvement on diagrams---but standard training does not
    route this signal across modalities.
  \item \textbf{Ignores asymmetric capability:} Per-view performance
    profiles differ across problems; a fixed training mix cannot
    leverage which view each problem is easiest from.
\end{itemize}

\section{Core Mechanism}

MIRROR introduces a per-problem view evaluation step before each RL
update. For each training problem $p$:

\begin{enumerate}[itemsep=2pt,parsep=0pt]
  \item Evaluate the current policy on all views
    $\mathcal{V} = \{T, I, T\!+\!I\}$ to obtain per-view rewards
    $r(p, v)$.
  \item Identify the best view: $v^*(p) = \arg\max_v r(p, v)$.
  \item Compute augmented rewards for weaker views that include a
    cross-modal alignment term encouraging the weaker-view distribution
    to match the best-view distribution.
\end{enumerate}

The policy update uses GRPO with the augmented rewards, simultaneously
training on all views. The reciprocal structure ensures that
improvement on the text view can feed into the image view and vice
versa.

\section{Technical Formulation}

Let $\pi_\theta$ be the policy and $a$ the correct answer.
The base reward is $r(p, v) = \mathbf{1}[\pi_\theta(p_v) = a]$.
The MIRROR-augmented reward for view $v \neq v^*(p)$ is:

\[
  \tilde{r}(p, v) = r(p, v) + \lambda
    \cdot \mathrm{KL}\!\left(\pi_\theta(\cdot \mid p_{v^*})\,
    \Big\|\, \pi_\theta(\cdot \mid p_v)\right)^{-1}
\]

The KL term penalizes divergence from the best-view response
distribution, rewarding weaker views for matching the best-view
reasoning trace. The inverse ensures the bonus is higher when the
weaker view is already close to the best view (facilitating transfer)
and smaller when the gap is large (avoiding overfit to a single view).

The full GRPO objective with MIRROR shaping:

\[
  \mathcal{L}_\text{MIRROR} = -\mathbb{E}_{p, v \sim \mathcal{D}}
    \left[\tilde{r}(p, v)\,
    \frac{\pi_\theta(a_v \mid p_v)}{\pi_{\theta_\text{old}}(a_v \mid p_v)}
    - \beta\,\mathrm{KL}(\pi_\theta \| \pi_\text{ref})\right]
\]

\section{Learning or Inference Procedure}

Training uses ODA-Data for all views simultaneously. Each mini-batch
contains all three views of the same problem, enabling within-batch
comparison of view-dependent rewards. The $\lambda$ coefficient
governing cross-modal shaping strength is annealed from 0.5 to 0.1
over training to prevent the alignment term from dominating late-stage
learning.

At inference, no change is needed: the model is queried with whichever
view is available. The improvement is baked into the model weights.

\section{What the Guarantee Says}

MIRROR does not guarantee cross-view performance parity---some problems
may remain fundamentally harder in one view than another. What it
provides is a systematic mechanism for reducing the gap, ensuring that
the model's best-view capabilities are not siloed and can transfer
across the view boundary.

\section{Experimental Findings}

\textbf{ODA-Data test set:}

\begin{center}\small
\begin{tabular}{lccc}
\toprule
Method & Text & Image & Combined \\
\midrule
LLaVA-1.6 (base) & 61.3 & 44.2 & 58.7 \\
+ GRPO only & 68.1 & 49.8 & 64.3 \\
+ MIRROR & \textbf{74.5} & \textbf{61.2} & \textbf{70.8} \\
\bottomrule
\end{tabular}
\end{center}

\textbf{External geometry benchmarks} (average over GeoQA, UniGeo,
Geometry3K): MIRROR $+$5.8pp over GRPO-only, $+$11.4pp over base.

The image-dominant view shows the largest absolute gain ($+$17pp over
base), consistent with the hypothesis that image reasoning benefits most
from cross-modal transfer from the textually stronger reasoning path.

\section{Ablations and Interpretation}

\textbf{Remove cross-modal shaping} ($\lambda = 0$): Image accuracy
drops 9.4pp relative to full MIRROR. The model still improves over
base from outcome RL, but the cross-view transfer is lost.

\textbf{Best view only} (train only on $v^*$, ignore weaker views):
Image accuracy drops 14.1pp relative to MIRROR. Focusing only on the
easiest view degrades weaker-view performance---the model over-specializes.

\textbf{Scrambled skills} (random KL target per problem): Image accuracy
falls to 47.1\%, below the GRPO-only baseline. The content of the
cross-modal alignment target---not just the fine-tuning signal---drives
the improvement.

\textbf{$\lambda$ annealing:} Fixed $\lambda = 0.5$ degrades image
accuracy by 3.2pp in late training, confirming that annealing is
important for preventing the alignment term from dominating the RL
signal.

\section*{Reference}

Wen Ye, Yuxiao Qu, Aviral Kumar, Xuezhe Ma.
\textbf{MIRROR: Learning from the Other View for Multi-Modal Reasoning}.
arXiv:2607.21552, July 23, 2026.
\url{https://arxiv.org/abs/2607.21552}

\end{document}
